% /root/monodromy, 2026-11-19. Proved; independent root and reduction_descent audits passed.
\subsection{A natural eight-dimensional block in every supersingular degree}
\label{subsec:ss-uniform-GL8-block}
\status{Proved and checked. The parent and reduction/descent agent
independently audited both boundary constructions, the full node-cover
comparison and the extraction of an actual natural block. This determines
a uniform subgroup. The full odd-degree images are determined later in
Theorem~\ref{thm:ss-all-odd-full-groups}, and the complete all-degree
classification is Theorem~\ref{thm:ss-product-all-full-groups}.}

\begin{theorem}\label{thm:ss-uniform-GL8-block}
For $X=E_0\times E_0$, $E_0:y^2+y=x^3$, with
$L=\mathcal O(3O)\boxtimes\mathcal O(3O)$ in characteristic two,
the degree-$d$ geometric Hasse--Witt image contains, for every $d\ge1$,
a subgroup acting as the natural $\GL_8(\mathbf F_2)$ on an
eight-dimensional subspace and fixing a complementary subspace pointwise.
Consequently the whole representation admits no invariant nondegenerate
alternating bilinear form. In even degree the same conclusion holds on
the common kernel of the invariant covectors of
Theorem~\ref{thm:ss-product-exact-covectors}.
No irreducibility or full-image assertion in higher degree is made here.
\end{theorem}

\begin{proof}
We work over an algebraic closure. The degree-one case is
Theorem~\ref{thm:ss-product-degree-one-full-group}.
We use the exact ranks of Theorem~\ref{thm:DE-A3-ss-product-rank}.
The maximum stable-rank open of the coherent family includes singular
divisors; its unit Lang system restricts to the local system in the
question on the dense smooth open, with the same image.
Indeed an iterated Frobenius image of constant maximum rank is a
unit subbundle, and restriction of a finite etale cover on a normal
integral scheme to a dense open preserves its connected components.

\paragraph{Odd degree: an odd self-isogeny and nodal images.}
Suppose $d\ge3$ is odd. Lemma~\ref{lem:DE-A3-ss-norm} gives an
endomorphism $b$ of $E_0$ with $b^\dagger b=bb^\dagger=[d-1]$.
Put
\[
 \phi=\begin{pmatrix}1&b\\-b^\dagger&1\end{pmatrix}:X\longrightarrow X.
\]
Then $\phi^\dagger\phi=\phi\phi^\dagger=[d]I$, and $\phi$ is
an etale isogeny of degree $d^2$ with constant odd kernel $G$.
For a generic smooth degree-one member $C\in|L|$, the image
$Z=\phi(C)$ is integral and nodal, and $C$ is its normalization.
Here is why the nodal-image argument applies also to this nonproduct
isogeny and to supersingular elliptic factors.

Only translations by points of the constant kernel enter that argument.
A generic $C$ is not invariant under any nonzero translation: if a
translation preserves $L$, its projective action is nontrivial because
$L$ is very ample; otherwise it preserves no member of $|L|$.
Thus $C\to Z$ is generically one-to-one. For a nonzero kernel point
$g$, double identifications are parametrized by
$C\cap(C-g)$, of degree $L^2=18$.

The tensor first-jet calculation of
Lemma~\ref{lem:product-small-degree-jets} uses only the plane cubic
embeddings and the fact that tangents through a fixed point are finite.
For the present cubic $Y^2Z+YZ^2+X^3$, the partial derivatives are
$X^2,Z^2,Y^2$, which are linearly independent, and no nonzero quadratic
vanishes on the cubic. Thus that finiteness condition holds.
Two first jets have rank at least five; rank six fails only on the
two coordinate-fibre loci and the finite crossed-tangency locus
over either fixed first point. For pairs $P,P+g$ with both coordinates
of $g$ nonzero, the exceptional points satisfy equations
$3P_i=\text{constant}$ and are finite. Away from them, the two value
conditions and proportionality of the two nonzero gradients impose
three conditions with $P$ varying in dimension two.
At the finitely many exceptions, the two independent values already
give a proper incidence image.

If one coordinate of $g$ is zero, the joint first-jet rank is five.
After imposing the two values, its sole remaining relation involves
the two derivatives along the translated factor and is nonzero;
the two transverse derivatives remain independent. The determinant
of the two gradients is not identically zero on this three-dimensional
allowed derivative space. Tangency is again a third condition.
Thus every $C\cap(C-g)$ is transverse for generic $C$.

Three-point identifications are excluded by the same Segre value
calculation used in Lemma~\ref{lem:ordinary-product-isogeny-nodal-images}.
Three distinct rank-one Segre tensors are collinear only if one factor
is constant and the other three points are collinear. For
$P,P+g,P+h$, the exceptional locus has dimension at most one, since
it requires translations in the same factor and
$3P_i+g_i+h_i=0$. The value ranks are three off this locus and two
on it; the incidence dimensions in $|L|=\mathbf P^8$ are at most
$8-3+2=7$ and $8-2+1=7$. Etaleness of $\phi$ therefore gives only
ordinary nodes. Each $g$ contributes an etale intersection cover of
degree eighteen, with the constant identification of $g$ and $-g$.
The node permutation image is consequently a subquotient of a
product of $S_{18}$'s.

The image has exactly the required line bundle. Write
$\Theta=(\{O\}\times E_0)+(E_0\times\{O\})$, so $L=\mathcal O(3\Theta)$.
The Rosati identity gives $\phi^*[\Theta]=d[\Theta]$.
For the norm line bundle, pulling back its class gives
\[
 \phi^*[\operatorname{Nm}_\phi L]
   =\sum_{g\in G}[t_g^*L]=d^2[L]=\phi^*[L^d].
\]
Pullback by an isogeny is injective on the torsion-free Neron--Severi
group after tensoring with $\mathbf Q$, so the two line bundles
have the same numerical class and differ by a member of $\Pic^0(X)$.
Both are symmetric, and $\Pic^0(X)[2](k)=0$; hence the difference
is trivial. Thus $Z\in|L^d|$ exactly.

Its arithmetic genus is $9d^2+1$, while the normalization has
genus ten and stable rank eight. Therefore it has $9(d^2-1)$
nodes and stable rank $8+9(d^2-1)=9d^2-1$, the full generic rank.
On this boundary family its Lang representation has an exact sequence
\[
                  0\longrightarrow T\longrightarrow V_d
                    \longrightarrow Q\longrightarrow0,
 \qquad Q\simeq V_1,\quad \dim Q=8,
\]
where $T$ is the graph permutation contribution. There is no ambient
stable contribution because $X$ is supersingular.
The quotient image is full $\GL_8(\mathbf F_2)$, whereas the whole
image on $T$ has prime factors at most seventeen.

\paragraph{Even degree: a fixed odd member and a symmetric node cover.}
Let $d\ge2$ be even. Choose a fixed smooth maximal-rank member
$D\in|L^{d-1}|$, which exists by the odd-degree rank theorem, and
vary $C\in|L|$ through its smooth maximal-rank open, transverse to $D$.
Then $Y=C\cup D$ is nodal and
\[
 f_Y=8+\bigl(9(d-1)^2-1\bigr)+18(d-1)-1=9d^2-3.
\]
This again equals the full generic rank. The boundary Lang system has
a filtration
\[
 0\subset T\subset T_0\subset V_d,
 \qquad T_0/T\simeq H^1_{\et}(D,\mathbf F_2)\text{ constant},
 \qquad V_d/T_0=Q\simeq V_1.
\]
The graph space is $T=\mathbf F_2^N/\mathbf F_2\mathbf 1$,
where $N=18(d-1)$ is the number of intersection points.
Its permutation image is the full $S_N$.

To justify that last assertion, the one-point incidence over $D$ and
the ordered distinct-pair incidence over $D\times D\setminus\Delta$
are irreducible projective bundles: the very ample embedding separates
one and two values. Thus the node action is doubly transitive.
There exists a section with exactly one double zero on $D$ and all
other zeros simple. Choose a general point $P\in D$ whose two
elliptic coordinates are not flex points. At $P$, $L$ separates
second jets, so the quadratic coefficient along the tangent of $D$
can be made nonzero after the value and first tangent derivative
are killed. For $Q\in D\setminus\{P\}$, the two curve jets at
$P,Q$ have combined rank four except for finitely many $Q$, and
rank at least three at those exceptions, by the preceding tensor
first-jet bounds. The exceptional points are finite because $D$ is
integral and ample, hence is not a coordinate fibre. Within the
codimension-two system imposed at $P$, another double zero imposes
two conditions with $Q$ varying in dimension one, or at least one
condition at finitely many exceptions. Its incidence image is proper.
The exact double zero is thus attainable without any second collision.
The point-incidence source is smooth and the base is smooth.
These are precisely the two conditions of
\cite[Proposition~3.1, arXiv version~3, pp.~7--8]{Entin2021}, already
used in Lemma~\ref{lem:all-X-independent-nodes}; they yield a
transposition also in characteristic two and therefore the full $S_N$.
Shrinking to the smooth maximal-rank transverse $C$ open does not
change this finite permutation image.

The action of $S_N$ on $\mathbf F_2^N/\mathbf F_2\mathbf1$ is
faithful for $N\ge3$: the difference of two distinct coordinate
vectors has support two and cannot be the all-ones vector.
The whole image on $T_0$ is an extension of $S_N$ by a two-group,
since $T_0/T$ is constant. Its only nonabelian composition factor
is $A_N$. This is not $\GL_8(\mathbf F_2)$: $11$ divides $|A_N|$
for $N\ge18$, whereas it does not divide
$|\GL_8(\mathbf F_2)|=2^{28}\prod_{i=1}^8(2^i-1)$.

\paragraph{Extraction of an actual block.}
In either parity let $T_0$ denote the kernel of the quotient to $Q$
(in odd degree take $T_0=T$). The simple group $\GL_8(\mathbf F_2)$
has no quotient in common with the entire $T_0$ image: in odd degree
use the prime $31$, and in even degree use the just determined
composition factors. Goursat's lemma shows that the subgroup acting
trivially on all of $T_0$ still surjects onto $\GL(Q)$.
Its translation kernel in $\operatorname{Hom}(Q,T_0)$ is
$Q^*\otimes T'$ for a subspace $T'\subset T_0$, since this module
is a direct sum of copies of the absolutely simple natural dual.
Modulo that kernel, the remaining block extension is a cocycle with
values in $\operatorname{Hom}(Q,T_0/T')$.
The proved elementary vanishing
$H^1(\GL_8(\mathbf F_2),Q^*)=0$
(Lemma~\ref{lem:char2-group-inputs}) makes it a coboundary.
After conjugation it is zero. The entire translation kernel is
present, so its elements then remove the remaining $T'$-valued
columns. Hence the subgroup contains every block diagonal matrix
$\operatorname{diag}(I_{T_0},A)$ with $A\in\GL(Q)$.
This proves the actual natural-block assertion.

Finally the natural eight-dimensional $\GL_8(\mathbf F_2)$ module
has no invariant covector and no nonzero invariant bilinear form.
For the latter, the stabilizer of any nonzero vector has no fixed
covector: its arbitrary upper row first kills the covector's first
coordinate, and its $\GL_7$ block kills the rest. Thus there is
no nonzero equivariant map $Q\to Q^*$.
Any invariant bilinear form on $V_d=Q\oplus T_0$ therefore vanishes
on $Q\times Q$ and on $Q\times T_0$, so $Q$ lies in its radical.
In even degree every invariant covector vanishes on $Q$; its common
kernel is $Q\oplus(T_0\cap K)$, and the same argument applies there.
This proves the stated form exclusions.
\end{proof}
